The judicial reasoning employed in \textit{Bradwell} reveals Component~3 ideological justification in its purest constitutional form. The concurring opinion authored by \textbf{Justice Joseph Bradley} codified the ideology of separate spheres into constitutional law with unusual candor: ``The natural and proper timidity and delicacy which belongs to the female sex evidently unfits it for many of the occupations of civil life. The constitution of the family organization, which is founded in the divine ordinance, as well as in the nature of things, indicates the domestic sphere as that which properly belongs to the domain and functions of womanhood.'' \textit{Bradwell}, 83 U.S.\ at 141 (Bradley, J., concurring). Bradley further specified that ``a woman had no legal existence separate from her husband, who was regarded as her head and representative in the social state.'' \textit{Id.}

This is \textbf{Component 3: Ideological justification} in its purest form. By framing structural exclusion as ``divine ordinance,'' the state naturalizes the asymmetry, recoding resistance as a challenge to a political algorithm, a status the divine framing had assigned to violations of the natural order. The same theological move Judge Bazile would perform in 1965 (``Almighty God created the races'')---a move Chief Justice Warren would explicitly name ``White Supremacy'' two years later in \textit{Loving}---was already fully operational in \textit{Bradwell} ninety years earlier, applied to the gendered partition.

\subsection{Protective Labor Laws and the Muller v. Oregon Paradox (1908)}

The state’s paternalistic justification for controlling women’s bodies and labor was further cemented during the Progressive Era through the implementation of ``protective'' labor laws. In \textit{Muller v.\ Oregon}, 208 U.S.\ 412 (1908) \cite{muller}, the Supreme Court unanimously upheld an Oregon statute that prohibited women from working more than ten hours a day in factories and laundries.

Justice Brewer, for the Court, grounded the decision in a biological theory of female incapacity that explicitly mimics the \textit{Bradwell} framework’s Component~3 logic: ``That woman’s physical structure and the performance of maternal functions place her at a disadvantage in the struggle for subsistence is obvious\ldots [A]s healthy mothers are essential to vigorous offspring, the physical well-being of woman becomes an object of public interest and care in order to preserve the strength and vigor of the race.'' \textit{Muller}, 208 U.S.\ at 421 (Brewer, J.). The ruling reinforced the concept that women were a distinct class requiring state guardianship---as minors do---by reason of their biological function.\footnote{Brewer explicitly invoked the parens patriae parallel: ``As minors, though not to the same extent, she has been looked upon in the courts as needing especial care that her rights may be preserved.'' \textit{Muller}, 208 U.S.\ at 421--422. The framework reads this as the state encoding a permanent biological proxy variable ($P_{\text{biology}}$) that is structurally identical to the racial proxy variables documented in Chapters~\ref{ch:enforcement} and~\ref{ch:kinetic}: an immutable physical characteristic is declared to produce legal incapacity, which justifies a restriction architecture managed by the state \textit{in the subject’s interest}.} Consequently, these protective laws combined limited relief from brutal factory conditions with severe restrictions on women’s earning potential and career mobility. Women were systematically fired from higher-paying positions---such as printers and streetcar conductors---and relegated to lower-wage, unregulated domestic and service sector jobs.

\subsection{The Intersectional Suffrage Paradox}

The narrative that the ratification of the Nineteenth Amendment on August 26, 1920, granted all American women the right to vote is a historical fallacy that obscures the systemic disenfranchisement of women of color. The amendment declared that voting rights could not be denied on the basis of sex; it left intact the vast architecture of state-sponsored voter suppression designed to exclude non-white populations.

For African American women, the victory of 1920 was largely illusory. State legislatures across the South and West utilized poll taxes, literacy tests, and unchecked domestic terrorism to prevent Black women from accessing the ballot box. This systemic exclusion persisted for nearly half a century until activists such as Fannie Lou Hamer and Diane Nash successfully championed the passage of the Voting Rights Act (VRA) of 1965.

Native American women were entirely excluded from the Nineteenth Amendment because they were not legally recognized as United States citizens in 1920. Activists like Zitkala-Sa, a Dakota woman who advocated for indigenous preservation and women's rights, tirelessly campaigned for Native enfranchisement. This advocacy contributed to the passage of the Indian Citizenship Act of 1924. Many local governments continued to employ literacy tests and poll taxes to block indigenous women from voting until successful litigation in New Mexico (1948), Utah (1957), and Nevada (1958).

Similarly, Asian American and Latina women confronted an overlapping matrix of exclusionary immigration and language policies. Beginning with the \textbf{Page Act of 1875}, which restricted the entry of Asian women under the guise of preventing ``immoral purposes,'' the federal government institutionalized the perception of Asian women as disease carriers and perpetual foreigners. The federal code still retains the scaffolding of this exclusion: \usclink{apx:usc:8:ch7}{8 U.S.C.~\S\S~262--297}---an entire chapter titled ``Exclusion of Chinese''---was not repealed until 1943:
\uscinline{usc_8_ch7}
The severe linguistic barriers preventing Asian American and Latina women from equitable voting access were only legally mitigated by the 1975 extension of the Voting Rights Act (Section 203).

\paragraph{Primary statutory text (VRA enforcement, Title~52).}
The modern codification of the Voting Rights Act's core vote-dilution standard appears at \usclink{apx:usc:52:10301}{52 U.S.C.~\S~10301}:
\uscinline{usc_52_10301}

\section{State Sanction of Intimate Violence and the Marital Rape Exemption}

The American legal system’s infringement on women's rights is illustrated by its historical tolerance---and outright legislative endorsement---of physical and sexual violence within the domestic sphere. The subjugation inherent in the doctrine of coverture extended seamlessly to a husband’s physical control over his wife's body. 

Perhaps the most egregious manifestation of state-sanctioned domestic violence was the legal exemption for \textbf{marital rape}. Grounded in the 1736 writings of English jurist \textbf{Sir Matthew Hale}, the marital rape exemption posited that the mutual matrimonial contract represented irrevocable, lifelong consent to sexual intercourse, meaning a husband could not legally be guilty of raping his wife. This archaic philosophy was maintained well into the modern era; the 1962 Model Penal Code explicitly defined rape as an act committed against a ``female not his wife.''

Criminalization reached all fifty states by 1993; many states nonetheless maintained lesser penalties or higher evidentiary thresholds for marital rape.

\begin{table}[htbp]
\centering
\small
\caption{Persistent Legal Disparities in Marital Rape Prosecution}
\label{tab:marital_rape}
\begin{tabular}{>{\raggedright}p{3cm} >{\raggedright\arraybackslash}p{9cm}}
\toprule
\textbf{State} & \textbf{Statutory Nuance} \\
\midrule
South Carolina & Requires force to be of a ``high and aggravated nature''; maximum penalty 10 years (vs 30 for non-spousal). 30-day reporting deadline. \\
Rhode Island & Excludes spouses from prosecution if the victim is ``physically helpless'' or mentally incapacitated. \\
California & Maintained separate criminal offenses until 2021 (AB 1171). \\
Virginia & Previously permitted therapy as an alternative to punishment to promote ``maintenance of the family unit.'' \\
\bottomrule
\end{tabular}
\end{table}

\section{The Breeding Apparatus: The Capitalized Womb and the Annihilation of Reproductive Autonomy}

The preceding sections have documented how European legal architecture stripped women of economic identity (coverture), civic participation (the ``divine sphere''), and physical sovereignty within marriage (the marital rape exemption). Under chattel slavery, these mechanisms converged into their most absolute form: the total commodification of the enslaved woman's reproductive capacity as a capital-generating asset for the Elite.

\subsection{The Legal Foundation: \textit{Partus Sequitur Ventrem} as Supply-Chain Engineering}

The 1662 Virginia statute of \textit{partus sequitur ventrem}---``the offspring follows the belly''---constituted a deliberate \textbf{reversal} of the English common law of patrilineal descent, in which a child's legal status followed the \textit{father}. Under English law, the children of an English father inherited his status regardless of the mother's condition. The Virginia legislature specifically inverted this principle to solve a supply-chain problem: white enslavers were raping enslaved women, and under patrilineal rules, the resulting children could claim free status through their fathers' lineage. By reassigning legal status to the maternal line, the Elite converted the systematic rape of enslaved women from a potential liability---children who might claim freedom---into a \textbf{production function}: every child born to an enslaved woman expanded the capital stock, regardless of paternity. The rapist's crime became the owner's profit.

This inversion simultaneously accomplished two structural objectives. First, it severed the enslaved child from any legal claim to the father's status, wealth, or lineage, ensuring that the offspring of interracial sexual violence remained permanently locked in $O_{\text{racialized}}$. Second, it transformed the enslaved woman's body into a \textbf{self-reproducing asset}: a machine that generated new capital units at no additional cost to the owner, incentivized by the very sexual violence the system authorized.

\subsection{The Breeding Apparatus After 1808}

The closure of the international slave trade in 1808 transformed this incentive structure from opportunistic exploitation into systematic industrial practice. With the external supply of enslaved bodies legally terminated, the domestic reproduction of the enslaved population became the \textit{sole mechanism} for expanding the capital stock. The enslaved woman's womb became, as historians Ned and Constance Sublette document, the ``capitalized womb''---the primary engine of wealth generation in the antebellum South \cite{sublette}.

The apparatus operated through overlapping mechanisms of coercion:

\begin{itemize}
    \item \textbf{Forced Pairings and ``Stock Men':} Enslavers exercised total control over the intimate lives of the enslaved, arranging unions without regard for existing bonds, preferences, or consent. Specific men were designated as ``stock men'' or ``studs'' and reserved---or hired out to neighboring plantations---for the sole purpose of impregnating enslaved women. Selection was based on physical strength, reducing human beings to livestock genetics. WPA Slave Narrative testimonies consistently describe this practice across multiple states.

    \item \textbf{Coerced Incestuous Rape and the Annihilation of Kinship:} The breeding apparatus did not stop at forced pairings between unrelated captives. The documentary record compiled by Sublette and Sublette, read alongside WPA Slave Narrative testimonies and the broader slave-breeding historiography, describes enslavers compelling incestuous rape---sons forced to rape mothers, brothers forced to rape sisters, fathers forced to rape daughters, and the reverse---under threat of the whip, the auction block, or the sale of kin \cite{sublette, baptist}. Thomas Foster's research on the sexual exploitation of enslaved men documents that boys and men were themselves coerced participants in this regime, their own bodies weaponized against the women they loved, and the women's bodies weaponized against them \cite{foster_rufus}. Structurally, this practice executed two subroutines simultaneously: \texttt{biological\_extraction.exe}, which under \textit{partus sequitur ventrem} converted the forced act into a new unit of capital regardless of paternity, and \texttt{kinship\_annihilation.exe}, which destroyed the pre-colonial African architectures of \textit{abusua}, patrilineal and matrilineal descent, and distributed intimacy documented in Chapter~\ref{ch:kinship} by turning kin into the instruments of one another's violation. No single act in the antebellum apparatus more completely collapsed the Out-group's remaining kinship immunity into the Elite's capital-production function. Black oral tradition has preserved this memory in the English term \textit{motherfucker}, which the African American vernacular received as a literal descriptor of the practice---a reading that historical linguistics, citing later written attestations, classifies as folk etymology even as it concedes that the word itself emerged from the descendant community.\footnote{The earliest confirmed written attestation is the 1889 Texas case \textit{Levy v.\ State}; subsequent usage is documented across the twentieth century. The absence of an antebellum paper trail does not disprove the oral-tradition origin; it reflects the documented fact that the archive of chattel slavery systematically suppressed records of the sexual crimes it authorized, producing precisely the asymmetry the framework predicts: the Elite's written record erases what the Out-group's oral record preserves. The oral tradition is itself a form of historical evidence---the descendant community's receipt for a crime the archive was structured not to record.}

    \item \textbf{Reproductive Regimentation:} Enslavers monitored menstrual cycles, pregnancy outcomes, and birth rates with the precision of agricultural yield tracking. Women who produced more children were rewarded with marginally better conditions; women who did not were subjected to punishment, forced re-pairing, or sale. Thomas Jefferson himself explicitly calculated the ``increase'' of his enslaved population as a primary financial consideration, treating birth rates as interest accruing on human capital \cite{baptist}.

    \item \textbf{Systematic Sexual Violence:} The breeding apparatus was inseparable from the regime of sexual terror documented throughout this chapter. Enslaved women were subjected to rape by enslavers, overseers, and their sons as a matter of routine. The threat of the whip enforced compliance with forced pairings. The Thistlewood diaries---recording 155 rapes of the enslaved woman Abba alone---are the operational logs of a system in which sexual violence was a production input.

    \item \textbf{Family Destruction as Economic Optimization:} Children born through forced reproduction were routinely separated from their mothers and sold, often before the age of ten, to meet the labor demands of the expanding cotton frontier. The domestic slave trade moved an estimated one million enslaved people from the Upper South to the Deep South between 1790 and 1860. Each sale was a capital transaction; each severed family was the system extracting the yield from the capitalized womb.
\end{itemize}

The scale of this apparatus defies euphemism. By the eve of the Civil War, the enslaved population of the United States had grown from approximately 700,000 at the time of the 1808 import ban to nearly \textbf{four million}---an increase generated almost entirely through forced domestic reproduction. The capital value of this population exceeded the combined value of all American railroads, factories, and banks. The enslaved woman's body occupied the center of the antebellum economy as the \textbf{primary engine of capital formation} in the wealthiest nation in the Western Hemisphere.

\subsection{The Colonial Specificity of the Breeding Apparatus}

The proto-partition framework established in Section~4.1 identifies this apparatus as a specifically European colonial innovation. In the pre-colonial West African societies from which the enslaved population was extracted, women exercised reproductive autonomy that the American system was specifically designed to annihilate. Akan women transmitted lineage and property through matrilineal descent. Igbo women's councils enforced community reproductive norms through collective authority. Yoruba women controlled their own fertility decisions within kinship structures that kept motherhood outside the capital production function.

The breeding apparatus applied the European gendered extraction template---the same architecture that produced coverture, the marital rape exemption, and the ``divine sphere''---to the intersection $O_{\text{racialized}} \cap O_{\text{gendered}}$, producing a compounding extraction rate that annihilated bodily autonomy along both axes simultaneously. The enslaved woman existed at the mathematical intersection of the two oldest partitions in the Elite's source code: she was property by race and reproductive machinery by sex. The breeding apparatus is the purest expression of what happens when both subroutines execute on the same body at the same time.

\section{Post-Mortem Extraction: The Ghost Value of the Enslaved Body}
\label{sec:post_mortem_extraction}

The preceding sections have tracked the extraction algorithm through the living body: the womb as a capital-production function under \textit{partus sequitur ventrem}, the breeding apparatus as a reproductive regimentation system, the coerced incestuous rape regime as simultaneous biological and kinship annihilation. Each of these subroutines operated on a subject who was, at minimum, still alive. The framework's internal logic, however, predicts a more complete result: if the enslaved body is fully converted to capital in the Elite's source code, there is no metaphysical reason for the extraction to terminate at biological death. The corpse, like the womb, should remain a production input. The historical record confirms this prediction precisely.

Historian Daina Ramey Berry's archival reconstruction of ``womb-to-grave'' pricing documents the phenomenon directly. Berry identifies a distinct category of valuation she terms \textbf{ghost value}: the monetary price enslaved bodies continued to command \textit{after} death, through the domestic cadaver trade, anatomical-specimen markets, dental prosthetics supply chains, and the commerce in preserved human material for everyday objects \cite{berry_ghost}. Ghost value is the empirical signature of a subroutine this framework names \texttt{material\_extraction.exe}: the conversion of the enslaved corpse (or living body harvested for non-reproductive biological yield) into raw material for Elite consumption. Where \texttt{biological\_extraction.exe} converts the womb into a capital factory, \texttt{material\_extraction.exe} converts the entire body---teeth, skin, hair, bones, organs---into commodity stock. The three subsections that follow document this subroutine at three phenotypically distinct material scales (teeth, skin, and hair); a fourth subsection states the Completeness Theorem for post-mortem extraction, and a fifth---Section~\ref{subsec:libidinal_extraction}---closes the broader taxonomy by documenting the in-vivo libidinal consumption of the living captive that sits alongside these material yields.

\subsection{Teeth: George Washington's Dentures and the Dental Prosthetics Supply Chain}
\label{subsec:teeth}

The single most rigorously documented instance of \texttt{material\_extraction.exe} in the American record sits inside the mouth of the nation's first President. On \textbf{May 8, 1784}, plantation manager Lund Washington recorded in George Washington's \textit{Ledger B, 1772--1793} (folio 179; Library of Congress) the entry: \textbf{``By Cash pd Negroes for 9 Teeth on Acct of Dr. Lemoire''}, at a price of £6 2s \cite{lund_ledger, mv_slave_teeth}. ``Dr. Lemoire'' was the French dentist Jean Pierre Le Mayeur, who corresponded extensively with Washington about dental work and about the procurement of human teeth for prostheses. Mount Vernon's own research historians, most notably Mary V.~Thompson, acknowledge the entry in full and situate it within the transatlantic commerce in human teeth that supplied the 18th- and early-19th-century dental profession \cite{thompson_unavoidable, mv_slave_teeth}. Whether those specific nine teeth were installed in Washington's own dentures is not forensically verifiable; his surviving dentures contain human teeth alongside animal teeth (cow, horse) and hippopotamus ivory, and the chain of custody for any individual tooth cannot be reconstructed. What is indisputable is the transaction: the President of the United States participated, directly and personally, in the commercial harvest of teeth from human beings his household legally owned.

The framework reading of this entry is strict, not rhetorical. Under chattel slavery, the legal capacity to refuse did not exist; payments recorded in a slaveholder's ledger to enslaved persons are transactions executed by the Elite upon captive bodies, not contracts between free agents. The ``payment'' is a bookkeeping artifact of the same system that recorded bushels of tobacco and head of cattle on adjacent pages. Thompson and the Mount Vernon historians are appropriately careful on the coercion question; the framework presented here resolves that carefulness by observing that coercion is the kernel condition of chattel slavery and requires no per-transaction proof. The broader supply chain was the era's \textbf{``Waterloo teeth''} market---so called because dentists sourced teeth from the battlefield dead after Waterloo (1815)---which extended seamlessly to the American plantation context, where the population of enslaved bodies constituted a stable domestic supply. Washington's ledger entry is a standard transaction within that market; it is a receipt.

The symbolic payload is equally exact. The founding executive of the United States used the bodies of people he held as property to form the words with which he governed a republic founded on the proposition that all men are created equal. The nation spoke, in its first decade, through the teeth of the enslaved. No further interpretation is required; the artifact is the argument.

\subsection{Skin: Nat Turner and the Anthropodermic Trade}
\label{subsec:skin}

The most extensively documented case of post-mortem skin extraction concerns \textbf{Nat Turner}, who led the 1831 Southampton County rebellion and was executed on November 11 of that year. The primary historiographical record is John W.~Cromwell's 1920 \textit{Journal of Negro History} article ``The Aftermath of Nat Turner's Insurrection,'' which states unambiguously that after the execution Turner's body was \textbf{dissected, skinned, his flesh rendered into grease, his skin fashioned into purses, and his bones distributed as trophies} \cite{cromwell_turner}. The extraction, however, was consumptive. Woodard recovers the testimony of William Sidney Drewry, who documented that Turner's flesh was rendered into grease and \textit{consumed} by the white populace as a grotesque form of anti-rebel medicine. This ritualized forced ingestion crossed the boundary from material harvesting into literal cannibalism, functioning as an ontological attempt to permanently absorb and neutralize the kinetic threat of the rebellion \cite{woodard_delectable}. Kenneth S.~Greenberg's edited volume \textit{Nat Turner: A Slave Rebellion in History and Memory} (Oxford University Press, 2003) treats these facts as historically established, as does the National Museum of African American History and Culture \cite{greenberg_turner, nmaahc_turner}. The material trail of Turner's remains persisted for more than a century: his skull circulated through private collections until it was returned to his descendants in 2016, a transaction reported in \textit{The New York Times} alongside the account of a contemporary student whose family retained what was described as a \textbf{purse made of Turner's skin} \cite{nyt_turner_purse}.

Turner's case is the named anchor of a systemic practice. Berry's research and Harriet A.~Washington's \textit{Medical Apartheid} jointly document the broader \textbf{domestic cadaver trade} through which Black bodies, living and dead, supplied the antebellum American medical establishment with dissection material, anatomical specimens, and, in specific recorded instances, preserved tissue for commercial and curio use \cite{berry_ghost, washington_medical_apartheid}. The Medical College of Georgia retained an enslaved man named \textbf{Grandison Harris} for decades as a grave-robber, dispatched specifically to rob Black cemeteries in Augusta for fresh cadaveric supply. The \textit{Mercury} newspaper of March 17, 1888 explicitly documented the fabrication of shoes, cigar cases, instrument cases, and book bindings from human skin and hair taken from cadavers in the American market \cite{jim_crow_museum_leather}. \textbf{Anthropodermic bibliopegy}---the binding of books in human skin---is documented in multiple 19th-century American specimens; many cannot be tied specifically to enslaved persons. The infrastructure that produced them (dissection rooms supplied by the domestic cadaver trade) was directly continuous with the extraction system under examination here.

The framework reading is once again strict. Turner rebelled against extraction and was converted, in death, into extraction. The state's reply to an act of anti-systemic violence was ontological: \textit{your body is raw material whether alive or dead, whether compliant or rebellious, and the algorithm does not distinguish between its outputs.} That the skin of the rebellion's leader became a purse carried, apparently, as a personal object into the twentieth century is the precise empirical confirmation of the framework's completeness thesis. The extraction is total.

\subsection{Hair: The North Georgia Chair and the Upholstery Economy}
\label{subsec:hair}

The third material scale---hair---occupies a different evidentiary register than teeth or skin. The primary piece of material evidence that has entered contemporary public consciousness is a \textbf{viral 2016 video}, reshared widely on social media in 2021, in which a furniture restorer working on a 200-year-old chair from ``a very wealthy family in North Georgia'' documented the removal of its upholstery and identified the interior stuffing as human hair \cite{newsweek_chair}. The restorer explicitly distinguished the material from the era's standard upholstery stuffings (horsehair, goat hair, pig hair), demonstrating on camera the texture, color, and curl pattern consistent with Black human hair and contrasting it with reference samples. Methodological honesty requires the observation that this artifact has not been subjected to forensic DNA analysis, and the chair itself remains in private hands. The restorer's visual identification is consistent with, but not strictly equivalent to, laboratory authentication.

The evidentiary gap limits authentication of specific objects; the broader practice is established by the documentary record. The 19th-century human hair trade is a historically documented supply chain, sourcing material for wigs, mourning jewelry, false hairpieces, and---as the 1888 \textit{Mercury} account cited above confirms---upholstery and cushion fill \cite{jim_crow_museum_leather}. Enslaved women's hair, in particular, was routinely cut or shaved as punishment, as a sanitary measure following illness, or as a byproduct of the sale process, and the resulting material was not returned to the persons from whom it was taken. The structural question shifts from whether human hair circulated as a commodity in antebellum America---it plainly did---to whether the specific biological source of upholstery stuffing in a given surviving object can be authenticated. In the case of the North Georgia chair, that authentication has not been performed; in the case of the broader practice, the documentary record is sufficient to establish its existence. The artifact is best read, then, as a named illustrative instance of an attested economic practice, with the broader practice established independently of this object's authentication.

The framework reading notes an important distinction between this subroutine and the previous two. The harvest of hair does not require the death of the source. It can be performed on the living body, repeatedly, over a lifetime, as part of the routine maintenance of the capital asset. Hair is therefore the \textbf{bridge case} between labor extraction (which operates on the living body during its productive years) and post-mortem material extraction (which operates on the corpse). Its inclusion here demonstrates that the \texttt{material\_extraction.exe} subroutine does not wait for biological death to execute; where a biological yield can be taken from the living captive without terminating the primary labor function, it is taken. The enslaved body is mined continuously across its entire temporal extent.

\subsection{The Completeness Theorem}
\label{subsec:completeness_theorem}

The three subsections above, read together with the preceding treatment of the breeding apparatus, yield a proposition this framework calls the \textbf{Completeness Theorem of the extraction algorithm}:

\begin{quote}
\textit{The algorithm's targeting of the Out-group body is temporally complete. Every phase of human biological existence---pre-conception, conception, gestation, birth, infancy, childhood, adolescence, adult labor, reproduction, old age, death, and post-mortem residue---is annexed by the algorithm as a production input. The enslaved body is capital at every scale and at every moment.}
\end{quote}

This completeness is what distinguishes chattel slavery from every other historically observed form of unfree labor. Roman slavery, as the Roman Control Case in Chapter~\ref{ch:portugal} has documented, was brutal and heritable under specific conditions but left the post-mortem body outside the extraction circuit; the corpse of a Roman slave did not enter the markets for dissection, rendering, tanning, or dental harvesting. Medieval European serfdom operated almost exclusively on the labor of the living. Even within the early modern Atlantic world, the completeness signature documented here---the conversion of the corpse into commodity stock, recorded in the private ledgers of the slaveholding class---is specific to the racialized extraction architecture developed in the Americas. The framework's thesis is that this completeness constitutes the signature of the specific ideological partition the Elite installed: the compilation step that reclassified enslaved Africans as property removed the legal and moral predicate that, in every other documented slavery regime, had at minimum preserved the corpse. Chattel slavery's totality is the signature of the ideological partition the Elite installed.

Formally, the Completeness Theorem's totality claim is supported by four named subroutines that together exhaust the temporal and material extent of the captive body: \texttt{biological\_extraction.exe} (the womb, \textit{partus sequitur ventrem}), \texttt{material\_extraction.exe} (the corpse and the bridge-case harvest of hair from the living source, documented in the three subsections above), \texttt{libidinal\_extraction.exe} (the in-vivo consumption of vitality and erotic value, documented at Section~\ref{subsec:libidinal_extraction}), and \texttt{scientific\_extraction.exe} (the experimental substrate, documented at Section~\ref{sec:medical_extraction}). No phase of the captive body's existence falls outside the union of these four subroutines; the algorithm's targeting is operationally complete.

\subsection{The In-Vivo Subroutine: \texttt{libidinal\_extraction.exe} and the Gastronomic Partition}
\label{subsec:libidinal_extraction}

The Completeness Theorem establishes that extraction is temporally complete. Between the post-mortem harvesting of \texttt{material\_extraction.exe} and the clinical exploitation of \texttt{scientific\_extraction.exe}, an evidentiary gap remains: the literal and libidinal consumption of the \textit{living} captive. Woodard's archive documents that the system actively consumed the Out-group's vitality to sustain the In-group. The extraction kernel operated a literal gastronomic partition. The theoretical ground for this reading is Orlando Patterson's comparative-sociological thesis in \textit{Slavery and Social Death}, which identifies slavery as a fundamentally \textit{parasitic} institution---one that the slaveholder camouflages, ideologically, by defining the slave as dependent when the parasitism runs in the opposite direction \cite{patterson_slavery, woodard_delectable}. Woodard's archival recovery demonstrates that Patterson's parasite metaphor was, in the nineteenth-century American record, a literal description.

First, Woodard recovers the suppressed archive of literal consumption. The records of the domestic trade, fugitive advertisements, and abolitionist testimonies are saturated with the rhetoric of consumption, bodily consumption, and hunger. The exemplary nineteenth-century naming of the kernel appears in John S.~Jacobs's 1861 \textit{A True Tale of Slavery}, which designates the American slaveholding class ``\textbf{human fleshmongers}''---vendors of human flesh, a term Jacobs applies with the literal inflection of the butcher's trade, stripping the word of any purely moral register \cite{jacobs_true_tale}. Solomon Northup's \textit{Twelve Years a Slave} describes the slaveholders' ``gastronomical enjoyments'' with the same literal inflection---a vocabulary of palate and appetite deployed on the body of the captive \cite{northup_twelve}. Harriet Jacobs's \textit{Incidents in the Life of a Slave Girl} renders the same culture of consumption operating inside the domestic space, with white mistresses acting out consumptive hungers in the intimate architecture of the plantation household \cite{jacobs_incidents}. The Works Progress Administration interviews of formerly enslaved persons, conducted in the 1930s, contain the literal-and-named anchor of the archive: the testimony ``\textbf{Master \ldots eated me when I was meat},'' which Woodard takes as the title-epigraph of his book and treats as the direct first-person record of the subroutine's operation \cite{woodard_delectable}. Collateral to the plantation archive, the 1838 voyage of the Nantucket whaleship \textit{Essex}---whose crew, stranded by cetacean attack, turned to cannibalism to survive---confirms the pattern at sea: the first four crew members murdered and eaten by the white survivors were Black, selected under a master/slave ideology that persisted even in a nominally Quaker-abolitionist vessel \cite{woodard_delectable}. The system produced what Woodard names a \textbf{Master Epicure} aesthetic---the conversion of extraction into a cultivated palate, in which the Elite's physical vitality was experienced as literally sustained by the ingestion of Black bodily energy.

Second, this extraction was fundamentally libidinal. The hunger of the slaveholding class was inextricably linked to sexual desire. The system extracted \textit{homoerotic and libidinal pleasure} through the absolute dominion over the captive body. The punishment rituals, the public auctions, and the intimate violence of the domestic space were structured as sites of libidinal consumption, establishing an enforcement subroutine ($F_{\text{enforce}}$) that operated through sexualized terror. The Master Epicure's appetite was both gastronomic and erotic---the same hand wielded the carving knife and the lash, and the same ideology coded the captive body as both meal and sexual object. The ideological source code for this double appetite is the Roman penetration hierarchy and its Augustinian theologization traced in Section~4.8: the ``buck breaking'' reference at the close of that section (``The American slaveholder who sexually violated an enslaved man was executing the same psychosexual logic'') is the line where Section~4.8's ideological genealogy and this subsection's archival record meet. Woodard's evidence is the empirical confirmation of a subroutine whose ideological specification the paper has already traced to Rome.

With the inclusion of this archive, the extraction taxonomy is closed. The four subroutines---\texttt{biological\_extraction.exe} (the womb), \texttt{material\_extraction.exe} (the corpse), \texttt{libidinal\_extraction.exe} (the living captive's vitality and erotic value), and \texttt{scientific\_extraction.exe} (the experimental substrate)---are jointly exhaustive. The algorithm targets the body at every scale.

\section{The Medical Extraction Laboratory: Scientific Progress on Black Bodies}
\label{sec:medical_extraction}

The Completeness Theorem established that the extraction algorithm annexes the Black body at every phase of biological existence. One subroutine within that framework demands separate treatment: the systematic use of Black bodies---living, captive, and non-consenting---as the raw material for scientific knowledge production. Where \texttt{material\_extraction.exe} converted the enslaved corpse into commodity stock, and \texttt{libidinal\_extraction.exe} consumed the living captive's vitality, \texttt{scientific\_extraction.exe} converted the living Black body into an experimental substrate, generating medical knowledge that accrued entirely to the Elite while inflicting costs---pain, disability, death, and intergenerational medical distrust---that remained entirely with the Out-group. This subroutine is the structural condition under which American gynecology, surgical technique, and infectious-disease research were founded and funded.

The subroutine's internal logic follows directly from the partition: once the Elite's source code has successfully reclassified a human population as property, the legal and ethical predicates that constrain experimentation on persons---informed consent, proportionate risk, compensation for harm---are automatically removed. What remains is an unregulated research substrate. American medical science operated within a legal architecture that had already made that exploitation structurally available, free of charge, and politically invisible.

\subsection{J.\ Marion Sims and the Founding of Modern Gynecology (1845--1849)}
\label{subsec:sims}

The physician designated the ``father of modern gynecology'' built his foundational surgical technique on the bodies of enslaved Black women in Montgomery, Alabama. Between 1845 and 1849, J.\ Marion Sims performed a series of experimental surgeries to repair vesico-vaginal fistulae---traumatic tears between the bladder and vagina caused by prolonged obstructed labor---on enslaved women whose names survive in part from his own published memoir: \textbf{Anarcha}, \textbf{Betsey}, and \textbf{Lucy} are the three most documented, with Anarcha alone subjected to at least \textbf{thirty surgical operations} before Sims successfully perfected his closure technique \cite{sims_story, washington_medical_apartheid, owens_medical_bondage}.

The anesthesia question requires precision. Diethyl ether had been demonstrated as a reliable surgical anesthetic at Massachusetts General Hospital in October 1846; chloroform followed in 1847. Sims's experimental operations span both the pre-anesthesia window (1845--1846) and the years when general anesthesia was clinically available (1847--1849). Contemporary scholarship, most rigorously in Deirdre Cooper Owens's \textit{Medical Bondage}, has complicated the flat claim that Sims performed all surgeries without anesthesia, noting that practice variations were common and that Sims's own memoir is internally inconsistent \cite{owens_medical_bondage}. The operative ideological framework is clear: the pseudoscientific claim running through Samuel Cartwright's taxonomy and much antebellum medical literature---that Black patients possessed physiologically diminished pain sensitivity---was factually baseless and ideologically functional \cite{washington_medical_apartheid, hoffman_pain_disparity}. Whether Sims administered anesthesia case by case or not, the architecture that made thirty operations on a single unanesthetized woman an acceptable experimental protocol was the same architecture that had classified her body as property: the partition removed the moral and legal cost of her suffering before any scalpel was raised.

The framework reading is strict. Anarcha, Betsey, and Lucy are the productive inputs of Sims's career. The surgical technique that Sims presented to the New York Academy of Medicine in 1852 and that founded the specialty of gynecological surgery was compiled on their bodies. That compilation produced professional prestige, financial reward, and scientific legacy for the Elite's knowledge-production apparatus; it produced pain, risk, and permanent bodily alteration for the Out-group without consent, without compensation, and with whatever anesthesia the operator chose to administer based on his own assessment of a racialized pain hierarchy.

The structural legacy is measurable in contemporary clinical data. A 2016 study in the \textit{Proceedings of the National Academy of Sciences} by Hoffman et al.\ found that a significant fraction of medical students and residents held false beliefs about biological racial differences in pain sensitivity---that Black patients have thicker skin, less sensitive nerve endings, or higher pain tolerance than white patients---and that these beliefs produced systematically lower and less accurate pain treatment recommendations for Black patients in clinical scenarios \cite{hoffman_pain_disparity}. The ideology Sims's era installed to justify experimental surgery on enslaved women is still executing in contemporary clinical settings, two centuries later. The subroutine persists; only the substrate has changed from the plantation operating table to the emergency room.

\subsection{The Tuskegee Syphilis Study: Withheld Cure as Algorithmic Precision (1932--1972)}
\label{subsec:tuskegee}

The Public Health Service study formally designated the ``Tuskegee Study of Untreated Syphilis in the Negro Male'' ran for \textbf{forty years}, from 1932 to 1972. It enrolled \textbf{399 Black men} in Macon County, Alabama who had latent syphilis, alongside 201 controls without the disease. Participants were promised free medical examinations, meals, and burial insurance. They were not told they had syphilis; they were told they were being treated for ``bad blood,'' a local idiom for various ailments \cite{jones_bad_blood, reverby_tuskegee}.

The study's structure was problematic from inception, but its moral collapse occurred in \textbf{1947}. In that year, penicillin became the standard of care for syphilis---a safe, proven, and curative treatment. The Public Health Service did not administer penicillin to the study participants. It \textit{actively intervened} to ensure participants did not receive it elsewhere: when Macon County men appeared on World War II military induction lists for standard syphilis treatment, PHS officials had their names removed from treatment rosters. The study continued for twenty-five more years. When PHS epidemiologist Peter Buxtun leaked the study's existence to the press in 1972 and the \textit{Washington Star} and \textit{New York Times} published exposés, the resulting Senate hearings produced the \textbf{National Research Act (1974)} and the Belmont Report's framework for human-subjects research ethics \cite{reverby_tuskegee}. By the time the study was terminated, 28 men had died directly of syphilis, 100 had died of related complications, 40 wives had been infected, and 19 children had been born with congenital syphilis.

The framework reading strips the study of the rhetorical insulation the word provides. The study was an active, multi-decade intervention to \textit{ensure} that a captured population of Black men remained exposed to a lethal pathogen after a cure existed---so that the Elite's knowledge-production apparatus could continue to extract data. The study's continuation after 1947 has no scientific justification independent of the implicit conclusion that the research value of the data exceeded the value of the participants' lives. That calculation is the algorithm's signature: the Out-group body as a production input whose welfare is an externality of Elite operation.

The aftershock is quantifiable. Economists Marcella Alsan and Marianne Wanamaker's 2018 analysis in \textit{The Quarterly Journal of Economics} found that in the years immediately following the 1972 revelations, Black men's engagement with the medical system declined measurably and persistently in counties with higher levels of Tuskegee awareness---producing a statistically significant increase in age-adjusted mortality among older Black men that the authors attribute directly to Tuskegee-driven medical distrust \cite{alsan_wanamaker}. The withheld cure killed in 1947. The knowledge that the cure was withheld then killed the next generation's willingness to seek care. The mortality differential compounded across decades. The extraction continues after a study ends; it schedules its own sequel through the population's embodied, behavioral, and institutional memory of betrayal.

\subsection{Operation Avon Park: Weaponized Entomology Over Black Communities (1955--1956)}
\label{subsec:avon_park}

Between 1955 and 1956, the United States Army Chemical Corps conducted covert open-air biological experiments over civilian populations in Florida. Under the Cold War biological-weapons research program, Army personnel released \textbf{millions of \textit{Aedes aegypti} mosquitoes}---the primary vector of yellow fever and dengue fever---over residential areas of \textbf{Avon Park, Florida}, a community whose residential geography was predominantly Black \cite{cole_clouds, washington_medical_apartheid}. The selection of target areas tracked racial geography. Populations less likely to generate political or legal backlash---and less likely to have their illness reports aggregated into a recognizable evidentiary pattern---minimized operational risk for the Army. The Out-group's structural political powerlessness was itself the targeting variable.

In the aftermath of the Avon Park releases, local residents reported illness clusters and deaths. Army personnel subsequently entered the affected area posing as public health workers, surveying residents and monitoring mosquito spread while concealing the operation's military origin \cite{cole_clouds}. The full scope of the domestic open-air biological testing program---encompassing hundreds of open-air experiments over American cities between 1949 and 1969, including releases of bacterial aerosols through New York City subway ventilation systems---was documented primarily through Freedom of Information Act litigation by researcher Leonard Cole, whose 1988 \textit{Clouds of Secrecy} remains the principal public account \cite{cole_clouds}. Senate investigations following the Church Committee exposures of the mid-1970s confirmed the program's scale and its largely covert domestic deployment against American civilian populations.

The framework reading of Avon Park extends the Completeness Theorem from the clinic to the residential neighborhood: where Sims used the operating table and the PHS used the diagnostic intake as the extraction mechanism, the Army Chemical Corps used the city block. The constant across all three is the designation of Black bodies as an acceptable experimental environment---a conclusion that requires, as its logical precondition, that the algorithm has successfully removed those bodies from the category of persons whose welfare generates a political cost to the Elite. The mosquito vector simply made the extraction airborne.

\subsection{The Radiation Archive and the Classified Biological Dossier}
\label{subsec:radiation_lyme}

The Cold War produced a parallel set of experiments involving ionizing radiation, documented through declassified federal records and most comprehensively catalogued by journalist Eileen Welsome in her Pulitzer Prize-winning investigation \textit{The Plutonium Files} (1999). Between the 1940s and 1970s, the U.S.\ Atomic Energy Commission and affiliated researchers administered plutonium injections, total-body irradiation, and other radioactive materials to hospital patients across the country, predominantly without their knowledge or consent \cite{welsome_plutonium, advisory_committee}. The experiments' purpose was military: to establish radiation tolerance thresholds in human subjects relevant to nuclear warfare.

The racialized selection logic is most explicitly documented at the University of Cincinnati, where Dr.\ Eugene Saenger---funded by the Department of Defense---conducted whole-body radiation experiments on \textbf{87 patients} between 1960 and 1971 \cite{advisory_committee}. The patient population was overwhelmingly \textbf{poor and Black}: 63 of 87 subjects were African American, and the median family income of participants was at or below the federal poverty line. The doses administered---ranging from 25 to 300 rads of total-body irradiation, with some patients receiving levels at which acute radiation syndrome and accelerated mortality occur---were not calibrated for oncological benefit. The patients were late-stage cancer patients, but the irradiation they received was calibrated to generate data for the Department of Defense, not to treat their disease. The Advisory Committee on Human Radiation Experiments, reporting in 1995, found that informed consent was not obtained in any form consistent with the minimal standards operative at the time \cite{advisory_committee}. Eight patients died within sixty days of their irradiation sessions.

A parallel classified biological-weapons research track implicates the \textbf{Plum Island Animal Disease Center} (New York) and \textbf{Fort Detrick} (Maryland) as potential origin points for ecological and epidemiological anomalies in the northeastern United States. Investigative journalist Kris Newby's 2019 \textit{Bitten}, drawing on deathbed interviews with Willy Burgdorfer---the Swiss-American entomologist who first isolated \textit{Borrelia burgdorferi} (the Lyme disease spirochete) in 1981 and who had spent decades conducting classified tick-based biological-weapons research for the Army using radioactively tagged arthropods as tracking vectors---presents circumstantial but documented evidence that experimental arthropod work during the 1950s and 1960s may have contributed to the ecological spread of tick-borne pathogens into the northeastern civilian population \cite{newby_bitten}. The claim remains contested and remains unconfirmed by declassified primary documentation. The 116th Congress passed \textbf{H.Res.~550} in 2019 directing the Department of Defense and the Department of Health and Human Services to investigate whether biological-weapons research contributed to the emergence of tick-borne illness in the continental United States \cite{hres550}. The congressional act of commissioning that investigation is itself evidence that the claim is regarded as sufficiently plausible by the nation's legislative body to warrant formal inquiry into whether the state's classified research apparatus deployed biological vectors against its own population.

The structural logic across the radiation experiments and the covert biological-weapons program is identical to Sims and Tuskegee: the Elite's knowledge-production and weapons-development apparatus required experimental inputs; the available inputs carrying the lowest political cost were those bodies and geographies whose structural position---poor, Black, institutionally dependent, politically unrepresented---rendered them least capable of refusal, complaint, or legal recourse. The algorithm produces a racially asymmetric outcome without explicit racial targeting. It requires only that the selection mechanism for ``acceptable risk'' operate on the same partition gradient the extraction architecture has maintained since 1662. Wherever that gradient exists, \texttt{scientific\_extraction.exe} will locate its substrate.

\section{Eugenics, Bodily Autonomy, and State Coercion}

Simultaneous to the economic marginalization of women, the state waged an aggressive legal campaign against female bodily autonomy. In 1873, Congress passed the \textbf{Comstock Act}, which classified contraceptives and abortifacients as ``obscene,'' making it a federal crime to transport such materials through the U.S. Postal Service. This legislative regime deliberately sought to keep women ignorant of reproductive control, ensuring that mandatory motherhood remained a defining characteristic of female existence. 

The state restricted the reproductive freedom of affluent women and simultaneously subjected marginalized women to violent reproductive coercion through the \textbf{eugenics movement}. In the 1927 case \textbf{Buck v. Bell}, the Supreme Court overwhelmingly upheld Virginia’s sterilization statute, determining that the state had a compelling interest in preventing the procreation of ``socially inadequate'' individuals. Justice \textbf{Oliver Wendell Holmes} famously wrote, ``Three generations of imbeciles are enough,'' solidifying the state's power to violently alter a woman's body against her will.

Following Buck v. Bell, more than 60,000 Americans were subjected to coercive sterilization, predominantly targeting Black, Latina, and Indigenous women. The systemic nature of this abuse was exposed in the 1970s cases of \textit{Relf v.\ Weinberger}, 372 F.\ Supp.\ 1196 (D.D.C.\ 1974) \cite{relf_v_weinberger}, in which Judge Gesell found that federally-funded sterilization without genuine informed consent was unlawful---the case arising from the non-consensual sterilization of Minnie Lee Relf (14) and Mary Alice Relf (12), two Black sisters in Montgomery, Alabama, whose illiterate mother signed consent forms she believed authorized birth-control injections---and \textit{Madrigal v.\ Quilligan}, No.\ CV 75-2057 (C.D.\ Cal.\ 1978) \cite{madrigal_v_quilligan}, in which Judge Curtis ruled against ten Mexican-American women who alleged coerced sterilization during labor at an Los Angeles county hospital, finding no proof of \textit{intentional} discrimination despite documented consent failures (the \textit{Washington v.\ Davis} intent-doctrine gate applied to a medical coercion context). This legacy persists: from the involuntary tubal ligations in California prisons (1997--2010) to the 2020 whistleblower reports of unwarranted hysterectomies in ICE detention centers in Georgia.

\section{Intersectional Violence: Lynching and Capital Punishment}

The violent enforcement of social hierarchies was exacted with singular brutality upon Black women. Between 1838 and 1969, an estimated \textbf{173 to 188 women were lynched} in the United States, the overwhelming majority of whom were Black women in the South. The lynching of Black women served a distinct sociological function: it was utilized as a weapon of racial terror to undermine Black families and mask an underlying epidemic of racialized sexual violence perpetrated by white men. 

The formal capital punishment system mirrors this history. Of the 356 documented executions of women between 1632 and 1962, 55\% were Black and 36\% were white. This racial disparity mirrors the broader hierarchy of the American algorithm, demonstrating that the Enforcement Class was weaponized primarily against women of color who deviated from the social order.

\section{The Modern Apparatus: Pregnancy Criminalization Post-Roe}

The overturning of \textit{Roe v.\ Wade}, 410 U.S.\ 113 (1973) \cite{roe_v_wade}, by \textit{Dobbs} in 2022 catalyzed an unprecedented acceleration in pregnancy criminalization. This modern apparatus bridges colonial ``concealing birth'' statutes with contemporary \textbf{fetal personhood} laws, effectively erasing the pregnant person's medical and legal autonomy. 

Cases like those of \textbf{Anne Bynum} (2015), \textbf{Latice Fisher} (2018), and \textbf{Brittany Watts} (2023) illustrate the transformation of healthcare providers into agents of law enforcement. Data compiled by \textbf{Pregnancy Justice} reveals that between June 2022 and June 2024, more than \textbf{412 women} were prosecuted after a birth, miscarriage, or abortion. In 399 of these cases, the primary allegation was substance use during pregnancy, often detected through non-consensual drug testing during medical care. This represents the final enclosure of the gendered axis: the legal reduction of the pregnant person to a hostile vessel, where the state's interest in the fetus supersedes the fundamental human rights of the woman.

\section{The Racialized Primal Wound: Why $O_{\text{racialized}} \cap O_{\text{gendered}}$ Produces Structurally Different Experiences}
\label{sec:racialized_primal_wound}

The preceding sections document how the gendered partition operates: coverture, civic erasure, marital rape exemption, the breeding apparatus, eugenics, and pregnancy criminalization. This chapter has treated $O_{\text{gendered}}$ as a unified category---women subjected to the extraction algorithm along the gender axis. But the intersection of $O_{\text{racialized}}$ and $O_{\text{gendered}}$ produces \textit{qualitatively different} oppression---different wounds, different relationships to the patriarchal contract, and different structural positions within the extraction architecture. This intersection is formalized in Chapter~\ref{ch:contradiction} (Equation~\ref{eq:12.8-intersection-double-partition}) as the set membership $O_{\text{racialized}} \cap O_{\text{gendered}}$, and its compounding effect is modeled by the multiplicative intersection equation (Equation~\ref{eq:12.9-multiplicative-intersection-compounding}): the effective extraction coefficient for a member of the intersection is $(1 - \alpha_r P_t) \cdot (1 - \alpha_g P_t)$, producing a multiplicative capacity reduction. The analysis that follows supplies the phenomenological and historical content of that formalism.

White women and Black women in America carry fundamentally different \textbf{primal wounds} from their respective histories within the gendered partition. Collapsing these into a single ``women's experience'' is analytically imprecise and obscures the structural divergence that shapes contemporary relational, political, and economic dynamics.

Throughout this section, ``primal wound'' names a \textbf{structural position} within the archive of law, labor, and kinship---an assigned coordinate in the extraction geometry. Wealth, region, queerness, disability, immigration status, and intra-racial class all modulate lived experience; the analysis isolates what tends to remain invariant \textit{across} those modulations when institutions and markets read the body through the same partition machinery.

\subsection{The White Feminine Wound: Enclosure Within the Contract}

White women's historical position within the gendered partition is one of \textit{enclosure within protection}. The ``wife script''---coverture, the domestic sphere, the obligation to function as an extension of a husband's legal and economic identity---was offered to white women as both cage and shelter. The wound is one of \textbf{overdetermined identity}: white women were \textit{seen} by the system, but only as wives, mothers, and domestic laborers. Personhood was submerged beneath function.

The feminist response to this wound is characterized by the pursuit of \textbf{autonomy as liberation}:

\begin{itemize}
    \item Freedom from the commitment script experienced as suffocating
    \item Bodily sovereignty as reclamation of a body that was legally subsumed under coverture
    \item Economic independence from a system that suspended legal identity upon marriage
    \item The dismantling of the ``divine sphere'' that Justice Bradley's \textit{Bradwell} concurrence codified
\end{itemize}

For white women, the patriarchal contract was \textit{offered and oppressive}. The logical response is to reject its terms---to seek a life outside the script.

\subsection{The Black Feminine Wound: Exclusion From the Contract}

Black women's historical position within the gendered partition is structurally inverted. The ``wife script'' was \textit{never offered}. Under chattel slavery, the legal architecture documented in Section~4.6---\textit{partus sequitur ventrem}, the breeding apparatus, the systematic rape of enslaved women as capital generation---reduced Black women to labor and reproductive capacity \textit{divorced from personhood}. Black women were excluded from the patriarchal contract entirely. The protections of coverture---however oppressive---at least acknowledged a woman as \textit{someone's} legal responsibility. Enslaved women received no such acknowledgment. Their bodies were property and their reproductive capacity was a capital asset; neither status derived from personhood or marriage.

The wound is therefore one of \textbf{devaluation}: value has historically been extracted from Black women's bodies without any corresponding recognition of personhood or entitlement to care. The wife script was never offered, so its potential burdens remain abstract. The lived reality is relational chaos, ambiguity, and the structural absence of guaranteed care.

The Black feminine response to this wound is characterized by the pursuit of \textbf{recognition as prize}:

\begin{itemize}
    \item Consistency, predictability, and promised care, received as acknowledgment of personhood; the same elements functioned as confinement for white women enclosed within the contract.
    \item Being seen as a \textit{person} with full standing, independent of any productive, reproductive, or sexual function.
    \item Non-transactional care: care independent of what one \textit{produces}
    \item The relational stability that white women experienced as a cage but Black women were denied entirely
\end{itemize}

For Black women, autonomy names the default state of being uncared for; white feminism celebrates the same condition as emancipation.

\subsection{The Structural Inversion and Its Contemporary Consequences}

This divergence has precise structural consequences:

\begin{enumerate}
    \item \textbf{Different relationships to feminist liberation}: The feminist project that liberated white women from the wife script may have made the script even \textit{less} accessible to Black women---universalizing an experience of autonomy-as-liberation that differs from an experience of autonomy-as-abandonment.
    
    \item \textbf{Different relationships to sex work}: White feminist perspectives tend to frame sex work as radical reclamation of bodily autonomy---taking control of one's market value on one's own terms. Black feminist perspectives tend to frame sex work as a continuation of the commodification that has operated since slavery: Black women's sexual and reproductive functions have \textit{always} been exploited outside of any social contract---owned in the objectifying, non-personhood sense, outside the patriarchal ``protective'' logic.
    
    \item \textbf{Different relationships to commitment}: When Black women seek consistency, exclusivity, and care from partners, this is the pursuit of a relational contract that was historically withheld. Conflating this pursuit with white feminist ``selective nostalgia'' for a system white women's foremothers inhabited is analytically imprecise and politically harmful.
    
    \item \textbf{The intersection compounds extraction}: Black men and Black women in the diaspora face a compounding problem. Black men are held to European monogamous standards within a system their ancestors did not choose (Section~3.2). Black women seek recognition within a relational structure from which they were excluded. Both are operating within a colonial imposition that serves neither's actual interests---and the conflict between them is another instantiation of the divide-and-conquer function, directing energy horizontally, away from $E$.
\end{enumerate}

\subsection{The Framework Implication}

The racialized primal wound reveals that the intersection $O_{\text{racialized}} \cap O_{\text{gendered}}$ is \textit{transformative}: the nature of the gendered wound changes depending on which side of the racial partition the woman occupies. Any analysis of the gendered axis that treats ``women'' as a monolithic Out-group will systematically mischaracterize the experiences of women at the intersection---either universalizing the white feminine wound (suffocation) and rendering the Black feminine wound (devaluation) invisible, or vice versa.

The extraction algorithm exploits this divergence. By producing \textit{different} wounds along the same axis, the system ensures that women cannot achieve cross-racial solidarity on gendered issues: they fight different fights because they were wounded differently. The white woman seeking escape from the contract and the Black woman seeking entry into the contract occupy structurally opposite positions. This is divide and conquer operating \textit{within} the gendered Out-group itself, along the racial axis. The Elite extracts from both.

\section{The Masculine Primal Wound: Why $O_{\text{racialized}} \cap I_{\text{masculine}}$ and $I \cap I_{\text{masculine}}$ Produce Structurally Inverted Experiences}
\label{sec:masculine_primal_wound}

The preceding section isolates the feminine primal wound across the racial partition. The framework's internal symmetry compels the corresponding analysis on the masculine axis. The intersection $O_{\text{racialized}} \cap I_{\text{masculine}}$ (Black men) and $I \cap I_{\text{masculine}}$ (white men) carry structurally inverted wounds within the same patriarchal contract, and those inversions generate their own divide-and-conquer geometry---both across the racial partition and, critically, \textit{within} the Out-group when the Black masculine wound collides with the Black feminine wound documented in Section~\ref{sec:racialized_primal_wound}.

\subsection{The White Masculine Wound: Enclosure \textit{as} the Contract}

White women were enclosed \textit{within} the patriarchal contract as its objects; white men were enclosed \textit{as} its enforcing subjects. They inherited the \textbf{provider-protector-patriarch script}---the reciprocal architecture of coverture: legal identity, property ownership, political sovereignty, the franchise, and the monopoly on organized violence (Section~6.1's dual policing genealogy is, among other things, the enforcement archive of this inheritance). Personhood was submerged beneath \textbf{function-as-enforcer}: head of household, breadwinner, sovereign citizen, armed defender of the partition.

The wound is therefore one of \textbf{overdetermined enforcement identity}---worth made contingent on the capacity to dominate, provide, and punish. Emotional life, vulnerability, and intimacy outside the performance of domination are structurally disallowed. The Roman penetration hierarchy and the Augustinian theologization of shame traced in Section~4.8 are the ideological source code: to be penetrated---economically, emotionally, sexually---is to fall out of the masculine In-group. The white man who feels the provider script as suffocating is inside a cage authored for him by the contract.

The masculine In-group response to this wound is characterized by \textbf{divestment from the performance}:

\begin{itemize}
    \item Freedom from the provider script experienced as suffocating---the cultural substrate of ``toxic masculinity'' discourse, ``soft masculinity,'' and the men's therapy wave.
    \item Emotional legibility as reclamation of an inner life that was legally and culturally conscripted into the sovereign role.
    \item Escape from the obligation to function as $F_{\text{enforce}}$---the armed, working, paying, punishing agent of the household and the state.
    \item Refusal of the ``divine sphere'' from the \textit{authoring} side: the same script Bradley wrote for women in \textit{Bradwell} wrote white men into the reciprocal cage.
\end{itemize}

For white men, as for white women, the patriarchal contract was \textit{authored and oppressive}. The logical response is to reject its terms---to seek a life outside the script.

\subsection{The Black Masculine Wound: Exclusion \textit{from} the Contract, Compounded by the Demand to Perform It}

The Black masculine position is structurally inverted. The provider-protector-patriarch script was \textbf{never granted}. The legal architecture documented in Chapters~\ref{ch:kinship}--\ref{ch:enforcement}---\textit{partus sequitur ventrem}, the prohibition on legal marriage under slavery, the prohibition on property ownership, the prohibition on arms (Chapter~\ref{ch:kinetic}'s gun-control genealogy), and the 13th~Amendment loophole---\textbf{excluded Black men from the patriarchal contract entirely}. They could not legally marry, protect, provide, vote, bear arms, or hold property. Their labor was property; their paternity was severed at auction; their capacity for violence was criminalized at the precise moment the white masculine capacity for violence was state-sanctioned through the slave patrols and their post-bellum successors.

Post-emancipation the exclusion persisted in converted form. Black men were then held \textit{to} the European monogamous, provider-protector standard they had been structurally barred from executing (Section~3.2's note on the imposed monogamous kinship template), while the material conditions for executing it---wealth, employment parity, physical safety, freedom from incarceration---remained systematically withheld. The lynching archive (Section~7.2) and the contemporary carceral archive (Chapter~\ref{ch:recompile}) are the enforcement record of this double-bind: Black men were punished with violence for \textit{successful} occupation of the role (Tulsa 1921, Rosewood 1923, the lynching of landowners and businessmen) and punished with pathologization for \textit{failed} occupation of it (``the deadbeat,'' ``the absent father,'' ``the thug'').

Mapped onto the sexual archive, the same inversion holds. White masculinity was scripted as the default-active penetrator of the Roman/Augustinian template; Black masculinity was scripted as \textbf{hypersexual threat}---the ``Black brute'' construction that exists precisely to justify the kinetic enforcement apparatus and the disarmament of Black men (Section~9). The same sexuality that granted white men sovereign standing was deployed against Black men as the rationale for lynching.

The wound is therefore one of \textbf{structural emasculation}---denial of the sovereign-enforcer role \textit{combined with} the demand to perform it, under pain of violence for either success or failure.

Against this structural emasculation, Woodard recovers the diasporic-reprise motif of the ``male daughter'' or ``male womb''---a cultural survival technology drawn from indigenous African social structures, specifically the Igbo male-daughter institution documented ethnographically at Section~\ref{subsec:male_daughter}. This was a radical mechanism of internal sustenance that recompiled an Igbo genealogical technology as a psychological survival strategy under captivity. To survive \texttt{libidinal\_extraction.exe}, the Black masculine subject developed an internal, self-nurturing capacity---a male womb---capable of generating the care and recognition the system structurally withheld \cite{woodard_delectable}. The Black masculine response to this wound is therefore characterized by \textbf{recognition as prize}---the exact mirror of the Black feminine response in Section~\ref{sec:racialized_primal_wound}:

\begin{itemize}
    \item Access to provider-protector standing as acknowledgment of personhood denied since arrival.
    \item Being seen as a \textit{man} in the legal-civic-sovereign sense, with standing independent of any labor, threat, or commodity function.
    \item The relational stability and legitimate kinetic capacity that white men experienced as a cage but Black men were structurally denied.
    \item The right to defend a household---the very function the enforcement apparatus has spent three centuries preventing.
\end{itemize}

For Black men, the ``divestment from masculinity'' that progressive discourse celebrates as emancipation is the continuation of a denial---one more demand that they abandon a sovereignty they were never granted in the first place.

\subsection{The Structural Inversion and Its Contemporary Consequences}

The four-item divergence that Section~\ref{sec:racialized_primal_wound} traces on the feminine axis runs symmetrically on the masculine axis:

\begin{enumerate}
    \item \textbf{Different relationships to ``liberation from masculinity''}: The progressive project that seeks to liberate white men \textit{from} the provider-protector script may render that script even \textit{less} accessible to Black men---universalizing an experience of masculinity-as-performance-burden that differs from an experience of masculinity-as-denied-status.
    
    \item \textbf{Different relationships to kinetic capacity and arms}: White masculine gun culture reads structurally as attachment to an inherited enforcement role. Black masculine pursuit of arms---the Deacons for Defense, the Black Panthers' open-carry patrols, the post-\textit{Bruen} surge in Black firearm ownership---reads as the Haitian Theorem (Chapter~\ref{ch:contradiction}) applied to the denied half of the contract: restoration of the kinetic variable from which the partition excluded them. Conflating the two is the same category error Section~\ref{sec:racialized_primal_wound} warns against on the feminine axis.
    
    \item \textbf{Different relationships to the provider role}: When Black men pursue head-of-household standing, exclusive partnership, and paternal authority, this is the pursuit of a civic-relational contract that was historically withheld. Coding it as ``toxic masculinity'' commits the same analytic error the paper warns against on the feminine side: universalizing the In-group wound and rendering the Out-group wound invisible.
    
    \item \textbf{The intersection compounds extraction}: Black men held to a monogamous provider standard their ancestors did not choose, and Black women seeking recognition within a relational contract from which they were excluded, are both operating inside a colonial imposition that serves neither---and the horizontal conflict between them is another instantiation of the divide-and-conquer function, precisely as the racial-axis conflict between white and Black women is on the feminine side.
\end{enumerate}

\subsubsection*{The Four-Quadrant Map of Primal Wounds}

Sections~\ref{sec:racialized_primal_wound} and~\ref{sec:masculine_primal_wound} together yield a complete $2 \times 2$ of primal wounds across the racial and gendered partitions:

\begin{center}
\begin{tcolorbox}[colback=black!3, colframe=black!70, arc=2pt, boxrule=0.6pt, width=\linewidth]
\renewcommand{\arraystretch}{1.35}
\small
\begin{tabular}{@{}p{0.14\linewidth}|p{0.40\linewidth}|p{0.40\linewidth}@{}}
 & \textbf{In-group (racial)}: $I$ & \textbf{Out-group (racial)}: $O_{\text{racialized}}$ \\
\hline
\textbf{In-group (masc.)}: $I_{\text{masculine}}$ & Enclosure \textit{as} the contract. Overdetermined enforcer-identity. Response: \textit{divestment from the performance.} & Exclusion \textit{from} the contract + demand to perform it. Structural emasculation. Response: \textit{recognition / access as prize.} \\
\hline
\textbf{Out-group (fem.)}: $O_{\text{gendered}}$ & Enclosure \textit{within} the contract. Overdetermined wife-identity. Response: \textit{autonomy as liberation.} & Exclusion \textit{from} the contract. Devaluation. Response: \textit{recognition as prize.} \\
\end{tabular}
\end{tcolorbox}
\end{center}

The diagonals of this matrix are the divide-and-conquer vectors the framework predicts and the historical record confirms: white women seeking \textit{out} of the contract, Black women seeking \textit{into} it (the racial-axis split of the gendered Out-group); white men seeking \textit{out} of the enforcer role, Black men seeking \textit{into} full sovereign standing (the racial-axis split of the masculine In-group). Four populations, same two axes, opposite vectors. This is the signature output of the Predatory Min-Max Function operating on an intersected partition---and it is why cross-quadrant solidarity has proven historically unachievable without the analytical reframe the framework supplies.

\subsection{Intersectional Out-group Contention: The Collision of the Black Masculine and Black Feminine Wounds}
\label{sec:intraracial_contention}

The framework's analytical payoff is sharpest where the two Out-group wounds meet inside a single relationship. Lay all three wounds the chapter has now isolated side by side:

\begin{itemize}
    \item \textbf{White masculine wound} ($I \cap I_{\text{masculine}}$): enclosure \textit{as} the contract $\rightarrow$ pursuit of \textbf{divestment from the enforcer role}.
    \item \textbf{Black masculine wound} ($O_{\text{racialized}} \cap I_{\text{masculine}}$): exclusion \textit{from} the contract + post-emancipation demand to perform it $\rightarrow$ pursuit of \textbf{recognition as sovereign / access to the provider-protector role}.
    \item \textbf{Black feminine wound} ($O_{\text{racialized}} \cap O_{\text{gendered}}$): exclusion \textit{from} the contract entirely $\rightarrow$ pursuit of \textbf{recognition as person, non-transactional care}.
\end{itemize}

The first two wounds diverge along the racial axis: one seeks to \textit{shed} the role, the other seeks to \textit{inhabit} it. The second two diverge along the gender axis \textit{inside the same racial Out-group}---and \textbf{this is the intraracial contention vector}. Both Black men and Black women are pursuing \textit{recognition}. But the relational postures through which each extracts recognition are inverted by the very patriarchal template neither of them authored.

\subsubsection*{The Structural Collision Points}

\begin{enumerate}
    \item \textbf{The recognition-posture inversion.} He pursues recognition by \textit{performing} sovereignty---headship, provision, protection, visible execution of a role historically denied. She pursues recognition by \textit{receiving} care---consistency, non-transactional presence, acknowledgment of personhood independent of function. The In-group's patriarchal template pairs these postures asymmetrically: his sovereignty is predicated on her subordination. Under that template, his recognition-seeking reads to her as \textit{demand for her subordination}; her recognition-seeking reads to him as \textit{withholding the sovereignty he is trying to access}. Each interprets the other's wound-response as personal injury; the imposition is structural.
    
    \item \textbf{The demographic double-bind.} Chapter~\ref{ch:kinship} documents that the male-biased transatlantic extraction collapsed African sex ratios and forced the universalization of polygyny---a structural adaptation to a demographic condition produced by the slave trade. The contemporary American carceral and employment apparatus \textit{reproduces the same demographic condition}: mass incarceration, disparate mortality, and disparate employment remove Black men from the partnership pool at rates that mirror the Middle Passage's extraction curve. Black men are then held to \textbf{European monogamous standards under demographic conditions that historically produced African polygynous structures}. The template and the material conditions are structurally incompatible. The resulting relational instability is read as moral failure on both sides; it is the predicted output of a colonial contradiction.
    
    \item \textbf{The role-performance trap.} His pursuit of the provider role requires material conditions---wealth accumulation, employment stability, property ownership, freedom from incarceration---that the Elite systematically withholds (Chapters~6 and~7). When he cannot execute the role, the shame is absorbed as personal failure even though the withholding is structural. She then bears the weight of a partner held to a standard he was architecturally prevented from meeting, and the frustration is displaced horizontally---onto him---and diverted from the apparatus that produces the withholding. The relationship becomes the execution environment for a conflict whose authorship lies entirely with $E$.
    
    \item \textbf{The care-asymmetry trap.} Her pursuit of non-transactional care requires a partner with the emotional bandwidth the enforcer-identity template leaves undeveloped. The Black masculine wound is \textit{compounded} by that template---the demand to perform sovereignty under hostile conditions leaves little residual capacity for the emotional legibility her wound requires. He arrives depleted by the performance demand; she arrives depleted by the absence of care. Each reads the other's depletion as indifference, when both are outputs of the same extraction.
    
    \item \textbf{The horizontal-energy conversion.} The intimate sphere becomes---precisely as Section~\ref{sec:racialized_primal_wound} predicts and the PHAM companion analysis documents---the site where the extraction algorithm converts potential Out-group solidarity into mutual exhaustion. Every hour spent litigating whether his provider-pursuit is patriarchal or her care-pursuit is conservative is an hour diverted from identifying that both wounds have the same author. The contention is the \textit{designed output} of two asymmetric wounds inflicted by the same partition, colliding inside an intimate sphere that was itself engineered (Section~4.4.1, nuclear-family atomization) to concentrate all relational stress onto a single dyad with no structural release valve.
\end{enumerate}

\subsubsection*{The Social-Reproduction Layer: When the Wound Trains the Next Node}

The contention extends beyond adult intimacy. It is reproduced generationally through household instruction, church culture, respectability politics, school discipline, dating scripts, and survival training. Black women are among its primary targets. But because the intimate sphere is also the first training environment for gendered behavior, Black women can become transmission nodes for patriarchal code under the very pressure misogynoir creates. A mother, aunt, grandmother, teacher, partner, or church elder may teach a Black boy that vulnerability is weakness, that emotional legibility is feminine failure, that heterosexual dominance is proof of manhood, that provision is the sole measure of worth, or that queerness and gender nonconformity are shame. The instruction is often misrecognized as protection: harden him before the world harms him. Structurally, it is the system converting a legitimate survival response into the next generation's enforcer script.

This social-reproduction layer is the overlooked bridge between misogynoir and Black male patriarchal performance. The system wounds Black women, then routes some of that wound into the manufacture of Black men trained to reproduce the gendered machinery that will later harm Black women, queer people, children, and themselves. The claim is diagnostic: a Black woman can be harmed by patriarchy and still participate in the social reproduction of patriarchal expectations. That participation harms her because it preserves the same partition that denied her care in the first place. Accountability for Black male harm remains necessary; alienation of Black men as a class is the interference engine's preferred conversion of that accountability into horizontal fracture.

\subsubsection*{The In-group Parallel and the Material Asymmetry}

The symmetrical contention exists on the In-group side: white men seeking divestment from the enforcer script and white women seeking autonomy from the wife script are also pulling in opposite directions along the same axis. But the In-group contention unfolds in a \textit{materially resourced} environment---the wealth, property, civic standing, and state protection the contract conferred remain largely intact for both parties even as they renegotiate their relationship to it.

The Out-group contention runs on the \textit{absence} of those resources. Same mechanism. Different material substrate. \textbf{Compounded harm.} This is why the same algorithmic collision produces materially different outcomes---marriage rate, divorce rate, intergenerational wealth transmission, carceral disruption of households---across the racial partition. The wounds are parallel; the material conditions for metabolizing the wounds are not.

\subsubsection*{The Framework Conclusion}

The Black-on-Black intimate sphere is one of the purest execution environments of the Predatory Min-Max Function the dataset contains. Two populations, each carrying an \textbf{inverted wound of the same colonial imposition}, are placed into a monogamous dyadic structure that concentrates all the friction, with no buffer resources and no analytical vocabulary---outside this framework---to identify the collision's author. The resulting contention is routinely coded as a \textit{cultural} problem (pathologized Black masculinity, pathologized Black femininity, ``the crisis of the Black family'') when it is, structurally, the precise relational output the extraction algorithm is engineered to produce.

The contention is the \textbf{signal} of the disease. Addressing it begins in the analytical reframe: naming the inversion, locating its author, and refusing the horizontal displacement that converts the wound into the weapon---so that the energy currently spent by Out-group men and Out-group women on each other can be redirected, for the first time since the partition was installed, at the partition itself.

\section{The Contemporary Intimate Sphere as Active Execution Environment}

The preceding sections trace the gendered partition's historical installation: from Roman penetration ideology through Augustinian theology, from coverture through the breeding apparatus, from the marital rape exemption through pregnancy criminalization. This historical trajectory might suggest that the gendered extraction kernel is a relic---a system whose legal architecture has been largely dismantled by feminist reform.

This suggestion is false. The extraction algorithm executes without legal architecture. It requires only \textbf{asymmetric autonomy distribution maintained through structural incentives}. And the contemporary intimate sphere---the bedroom, the household, the relationship---is a live execution environment where the algorithm runs in real time.

A companion analysis, \textit{Pseudo-Hybrid Asymmetric Monogamy} (PHAM), documents this contemporary execution in detail. Using the same Predatory Min-Max Function and 5-tier operational hierarchy defined in this framework, the PHAM analysis demonstrates:

\begin{itemize}
    \item \textbf{The intimate sphere reproduces the extraction architecture}: One partner retains complete autonomy (bodily, relational, economic, sexual) while exercising authority over the other partner's autonomy---the same Predatory Min-Max \textit{logic} instantiated at dyadic scale. \textbf{Mapping discipline}: both partners remain outside the Elite class ($E$); the macro five-tier hierarchy remains in force off the dyad (capital, cultural interface, institutional enforcement). Within the couple, advantaged versus disadvantaged positions align with the In-group versus Out-group extraction relation characteristic of the buffer architecture---without elevating either partner to ruling-class status.
    
    \item \textbf{The nuclear family's atomization (Section~4.4.1) creates the scarcity that intimate extraction exploits}: Monogamy concentrates all intimacy needs on a single partner, creating artificial scarcity. When that partner controls sexual and emotional access while enforcing exclusivity, the disadvantaged partner is enclosed---all legitimate outlets are closed.
    
    \item \textbf{The gendered partition's divide-and-conquer function executes through intimate conflict}: Partners depleted by relational asymmetry and the double burden of dual-income households lack the energy for political consciousness or collective action. The intimate sphere becomes a site where the extraction algorithm converts potential solidarity into mutual exhaustion---the same mechanism documented on the racial axis in Chapter~\ref{ch:containment}.
    
    \item \textbf{The commodification of intimacy feeds capital}: Dysfunctional relationships drive consumption---therapy, self-help, retail therapy, separate households after divorce---each a revenue stream for $E$. The system profits from the dysfunction it produces, just as the carceral system profits from the crime it manufactures.
\end{itemize}

The framework's implication is that the gendered partition documented in this chapter is \textbf{actively executing} through the structural incentives, cultural programming, and atomized relational architecture that feminism left intact. The algorithm adapted. It recompiled. And it continues to extract.

\begin{keyinsight}[Hardware Reading: The P-N Junction of Gender]
The gendered axis operates as the second imaginary component of the Complex Wage: $j_2 \psi_{s,\text{gender}}$. In circuit terms, coverture was a \textbf{P-N junction} installed across the reproductive current: it permitted conduction in one direction only---male inheritance, male property rights, male legal personhood---while blocking reverse current (female autonomy, female property, female citizenship). The 1973 \textit{Roe} decision briefly lowered the junction's forward voltage, allowing some reverse current; \textit{Dobbs} (2022) raised the breakdown voltage and restored unidirectional conduction. Eugenics (\textit{Buck v.\ Bell}, 1927) was a \textbf{Zener diode} in the reproductive circuit: when the Out-group's fertility threatened to alter the demographic charge distribution, the state authorized forced sterilization to clamp the voltage. The gendered partition is a parallel branch of the same extraction topology, drawing current from the same $E$ source and routing it to the same $O$ ground.
\end{keyinsight}
